% !PRINTED 405
\weberpage{405}
der Entscheidung unserer Frage $\nu\leqq\frac{1}{2}n$ annehmen. Ausserdem
folgt aus \S.\,2, dass die Coëfficienten $\alpha$ und $\beta$ von $\varphi$ und $\psi$
\gesperrt{ganze Zahlen} sein müssen.

Um nun über die Möglichkeit einer solchen Theilung zu
entscheiden, nehme man die Function $\varphi(x)$ mit unbestimmten
Coëfficienten $\alpha$, und dividire $f(x)$ durch $\varphi(x)$ (nach \S.\,3). Es
ergiebt sich dann ein Rest vom Grade $\nu-1$, und wenn man
dessen Coëfficienten, die sämmtlich rationale Functionen der $\alpha$
sind, gleich Null setzt, so erhält man $\nu$ Gleichungen, denen diese
Coëfficienten genügen müssen. Man kann durch Elimination (\S.\,50)
eine Gleichung herstellen, die nur noch eine dieser Unbekannten
enthält, und die dann eine ganzzahlige Wurzel haben muss. Es
wird aber oft zweckmässiger sein, diese Elimination nicht wirklich
auszuführen,. sondern direct zu versuchen, dem System der % [sic] stray comma-period as printed
$\nu$ Gleichungen durch ganzzahlige Werthe der $\alpha$ zu genügen. Die
Auswahl der zu erprobenden Zahlen kann durch manchen Kunstgriff,
der auf zahlentheoretischen Sätzen beruht, sehr eingeschränkt
werden.

Wir wollen beispielsweise eine Gleichung zehnten Grades
aus der Theorie der elliptischen Functionen entnehmen, wo
solche Aufgaben sehr häufig in der Weise vorkommen, dass die
Zerlegbarkeit in Factoren bestimmten Grades theoretisch feststeht,
und wo es sich dann darum handelt, diese Factoren zu
finden $^1$). Eine solche Gleichung ist
\[
x^{10} - 36\,x^8 + 528\,x^6 - 3897\,x^4 + 14354\,x^2 - 21025 = 0.
\]
Wir fragen, ob die linke Seite in zwei Factoren fünften
Grades zerlegbar ist. Wenn der eine der beiden Factoren
\[
x^5 + \alpha x^4 + \beta x^3 + \gamma x^2 + \delta x + \varepsilon
\]
ist, so ist der andere, da in der gegebenen Gleichung keine
ungeraden Potenzen vorkommen,
\[
x^5 - \alpha x^4 + \beta x^3 - \gamma x^2 + \delta x - \varepsilon,
\]
und es muss also sein:
\[
\begin{aligned}
x^{10} - 36\,x^8 + 528\,x^6 - 3897\,x^4 + 14354\,x^2 - 21025 \\
= (x^5 + \beta x^3 + \delta x)^2 - (\alpha x^4 + \gamma x^2 + \varepsilon)^2.
\end{aligned}
\]

\noindent\rule{4cm}{0.4pt}

$^1$) Vgl. des Verfassers Werk: „Elliptische Functionen und algebraische
Zahlen“, Braunschweig 1891. Das hier behandelte Beispiel aus der Theorie
der Transformation 47$^{\text{ten}}$ Grades ist aus der Abhandlung genommen: „Ein
Beitrag zur Transformationstheorie der elliptischen Functionen etc.“ von
\gesperrt{H.\ Weber} (1893). Mathematische Annalen, Bd.\ 43.
